\section{Evaluation Design}
\label{sec:methodology}

\subsection{Research questions and comparison target}
\textbf{RQ1:} How do project-context workflows change independently verified
completion of development tasks, and what resources do they consume?

\textbf{RQ2:} How much preparation, checking and repair work is required to keep
context accurate as the same project changes?

The primary outcome is verified task completion under a common execution-limit
policy. Resources and upkeep are separate secondary outcomes. These are design
choices for the research; they do not configure or execute the experiments.
Section~\ref{sec:comparison} defines the pairwise decision rule. The study tests
complete workflows, so unequal information is permitted and documented.

\begin{table}[htbp]
\centering\small
\caption{Methods and the contrast each permits.}
\label{tab:methods}
\begin{tabularx}{\textwidth}{@{}>{\raggedright\arraybackslash}p{0.15\textwidth}Y Y@{}}
\toprule
Method & Available context & Interpretation \\
\midrule
A: existing workflow & Native Memos instructions, documentation and source through the selected agent's normal tools. & Reference practice; already contains context. \\
\addlinespace
B: prepared context & The same native material plus deliberately prepared context, frozen before scored work. & Incremental value of a documented preparation workflow. \\
\addlinespace
C: proposed strategy & Undesigned; retained as a shell. & No scored contrast until its inputs, preparation and delivery are specified. \\
\bottomrule
\end{tabularx}
\end{table}

B's eventual specification must identify the preparation sources, author/tool,
review procedure, loading rules and update policy. Preparation may use source
and public documentation, but not scored task solutions or hidden evaluation.
Task-specific preparation, if later permitted, must be declared symmetrically
as an available workflow stage and fully costed. Revisions learned from scored
outcomes belong to a later study, not the original comparison. C is not assumed
to be retrieval-based.

\subsection{Pinned inputs, isolation and evaluation flow}
The candidate source revision is
\nolinkurl{b40d8a7b573774c52b022c8d6537d6baf2129dc0}. It was inspected for
project documentation; it is not an experimentally validated baseline.
Before scoring, verify the build environment and baseline checks, freeze the
revision, dependencies, task statements, acceptance criteria and evaluator,
and record their hashes. All methods receive the same task and starting code.

For each task and repetition, use a fresh conversation and independent writable
checkout. Apply the method's frozen context, run the agent, close the session,
seal the final patch and usage record, then evaluate the patch independently.
The evaluator runs on a fresh copy of the same baseline. It must not trust the
agent's declaration of success or its edited tests as the sole oracle.

Keep model/version, agent configuration, tool access, network policy, environment
and execution-limit policy common within the comparison. The limit policy must
state exactly which of tokens, time and tool steps are capped and how termination
works; equal limits do not imply equal actual consumption. Randomize method
order within task/repetition blocks, balancing scheduling across methods.
Record provider drift and any reproducibility controls that are unavailable.

Separate access permissions, processes, conversations, caches, writable homes
and persistent memory as applicable. Folder duplication alone does not enforce
isolation. Neither coding agents nor context preparers may read hidden checks,
reference patches, another method's context, or its outputs. Record actual
context loading: available text is not evidence that an agent consumed it.

\subsection{Task selection and independent acceptance}
Define the intended Memos workload and task eligibility before observing
comparative success. Sample across the chosen change surfaces rather than
selecting tasks because prepared context already contains their answers.
Record provenance, public-solution exposure and exclusions. A convenience task
set supports a conclusion about that set; broader inference requires a justified
sampling argument. No tasks have yet been selected.

Use separate pilot tasks to validate installation, task clarity, acceptance
coverage and analysis feasibility. Do not include preparation-tuning tasks in
scored evaluation. Validate checks against the unchanged baseline and a known
acceptable implementation where available. Behavioral bug checks should fail
before a repair and pass after it; existing behavior checks should already pass.
Review test sensitivity with deliberately invalid changes when feasible.
Any convention or design requirement needs an explicit executable check or a
frozen review rubric, with reviewers blinded to the method where possible.

This follows the independent issue-to-tests principle used by
\citet{jimenez2024}. Passing checks establishes the measured acceptance criteria,
not exhaustive software quality. Keep the evaluator identical across methods.
If evaluation is found defective, preserve the initial verdicts and apply a
method-blind correction to every affected artifact; document the deviation.

\subsection{Outcome and resource records}
For method $m$, task $t$ and repetition $r$, define $Y_{mtr}=1$ only when the
final artifact passes all frozen task-acceptance checks, required regression
checks and change-permission checks within the execution policy. Otherwise
$Y_{mtr}=0$. Record individual check outcomes and regression failures as well as
this binary endpoint. An agent-added test, a plausible patch or a loaded context
file is not sufficient for success.

Timeouts, invalid patches, agent-caused errors and unauthorized manual rescue
count as failures. Before scoring, define infrastructure-failure criteria and a
balanced rerun policy independent of the observed winner. Preserve original
attempts and their consumption; do not silently drop failures or retain only
the best retry. Unknown verdicts must be shown separately with best/worst-case
sensitivity bounds. If missing outcomes can change a claimed winner, the
conclusion remains inconclusive.

\begin{table}[htbp]
\centering\small
\caption{Required records and their accounting boundaries.}
\label{tab:measurements}
\begin{tabularx}{\textwidth}{@{}p{0.20\textwidth}Y@{}}
\toprule
Record & Definition and boundary \\
\midrule
Correctness & Task/repetition identifiers, acceptance and regression verdicts, failure category, baseline/evaluator hashes and final artifact. \\
\addlinespace
Tokens & Provider-reported input, cached input, output and reasoning categories with their overlap semantics; all attempts, including preparation calls. \\
\addlinespace
Elapsed time & Separate preparation, execution, updating and evaluation intervals; include execution timeouts and retries. \\
\addlinespace
Human work & Observed active preparation, review, intervention and maintenance minutes; concurrent waiting is not active work. \\
\addlinespace
Context provenance & Context version, permitted sources, factual review, preparation procedure and evidence of actual loading. \\
\addlinespace
Maintenance & Shared source-change identifier, audited claims, defects, repairs, unresolved cases and observed update/review resources. \\
\bottomrule
\end{tabularx}
\end{table}

Cached tokens may be part of input and reasoning tokens part of output: do not
add overlapping categories twice. Document size is not consumed usage. Missing
usage remains unknown, not zero. Monetary cost requires actual charges or dated
rates mapped to compatible categories. Report evaluator/research overhead
separately from costs an adopting team would incur.

For a declared workload of $N$ executions and resource component $j$, report
\[
 R_{m,j}(N)=P_{m,j}+\sum_{i=1}^{N}E_{m,i,j}+\sum_k U_{m,k,j},
\]
where $P$ is initial preparation, $E$ execution and $U$ context updating/review.
Apply the equation separately to compatible units. Include failures and retries.
Show initial and recurring amounts and the observed reuse horizon; extrapolated
amortization is a labelled scenario, not observed company savings.

\subsection{Maintenance episodes}
RQ2 needs source evolution, beyond the static task comparison. Supply the same
preselected code changes to each method's own context history, retaining no
cross-method learning. Label changes as natural or designed and record their
selection. Audit relevant guidance immediately after each change, then again
after the method's declared update procedure. Keep task scoring at matched
source revisions separate from these maintenance observations.

Classify audited claims as supported, stale, unsupported or unresolved, retaining
source evidence and scope. Report inspected and affected denominators, known
error rates, unresolved fractions and active checking/repair time. Empty or
incomplete context must not earn an accuracy advantage merely by making fewer
claims: separately audit coverage of a shared set of relevant requirements.
Different scope can explain divergent instructions; divergence alone is not a
contradiction. Edit frequency does not measure factual accuracy or human labor.

A maintenance-efficiency claim requires comparable post-update accuracy and
coverage as well as lower effort. Prespecify any acceptable degradation margins
before seeing outcomes. Keep operator experience and learning/order effects
visible and balanced where feasible. Episodes within one project trajectory
are dependent; report that trajectory as a case study unless independent
trajectories support formal inference. With no observed episodes, RQ2 remains
unanswered.
